\documentclass[12pt, a4paper, twoside]{book}

% --- Lingua e Codifica ---
\usepackage[italian]{babel}
\usepackage{iftex}

\ifPDFTeX
  % Per compilazione standard pdfLaTeX
  \usepackage[T1]{fontenc}
  \usepackage[utf8]{inputenc}
  \usepackage{lmodern}
\else
  % Per compilazione XeLaTeX o LuaLaTeX (Consigliato su Overleaf)
  \usepackage{fontspec}
  \usepackage{unicode-math}
  \setmainfont{DejaVu Serif}
  \setsansfont{DejaVu Sans}
  \setmonofont{DejaVu Sans Mono}
  \setmathfont{Latin Modern Math}
\fi

% --- Geometria e Layout ---
\usepackage[
  a4paper,
  twoside,
  inner=3.5cm,
  outer=2.5cm,
  top=2.5cm,
  bottom=2.5cm
]{geometry}

\usepackage{setspace}
\usepackage{emptypage}

% --- Grafica e Tabelle ---
\usepackage{graphicx}
\usepackage{caption}
\captionsetup{font={small, it}}

% --- Matematica ---
\usepackage{amsmath}
\usepackage{amssymb}

% --- Configurazione Intestazioni e Piè di pagina ---
\usepackage{fancyhdr}
\pagestyle{fancy}
\fancyhf{} 
\fancyfoot[LE,RO]{\thepage} 
\renewcommand{\headrulewidth}{0pt}
\renewcommand{\footrulewidth}{0pt}

% Ridefinisce lo stile 'plain' (per le pagine d'inizio capitolo)
\fancypagestyle{plain}{%
    \fancyhf{}
    \fancyfoot[LE,RO]{\thepage}
    \renewcommand{\headrulewidth}{0pt}
    \renewcommand{\footrulewidth}{0pt}
}

% --- Formattazione Titoli ---
\usepackage{titlesec}
\titleformat{\chapter}[hang]%
  {\normalfont\Huge\bfseries}% 
  {\thechapter.}%
  {1em}% 
  {}% 

\titlespacing*{\chapter}{0pt}{50pt}{40pt}

% --- Bibliografia ---
\usepackage[numbers]{natbib}

% --- Link ipertestuali (Caricare per ultimo) ---
\usepackage{hyperref}
\hypersetup{
    colorlinks=true,
    linkcolor=blue,
    citecolor=blue,
    filecolor=magenta,      
    urlcolor=cyan,
    pdftitle={Tesi Intelligenza Artificiale - Simone Garau},
    pdfauthor={Simone Garau},
}

% --- Documento ---
\begin{document}

% --- Frontespizio ---
\frontmatter

\begin{titlepage}
    \centering
    % Assicurati che il percorso dell'immagine sia corretto su Overleaf
    
\includegraphics[width=8cm]{img/logo_upo.jpg} 
    \vspace{1cm}
    
    {\large Università degli Studi del Piemonte Orientale \par}
    {\large Dipartimento di Scienze e Innovazione Tecnologica \par}
    \vspace{0.5cm}
    {\large \bfseries Corso di Laurea Magistrale in Intelligenza Artificiale e Innovazione Digitale \par}
    
    \vfill
    {\Large Tesi di Laurea Magistrale \par} 
    \vspace{1cm}
    {\huge\bfseries Tesi Garau \par}
    \vspace{0.5cm}
    {\Large embedding e XAI \par}
    \vfill

    \begin{minipage}[t]{0.45\textwidth}
        \begin{flushleft}
            {\bfseries Relatore:}\\[0.2cm]
            Prof. Giorgio Leonardi
        \end{flushleft}
    \end{minipage}
    \hfill
    \begin{minipage}[t]{0.45\textwidth}
        \begin{flushright}
            {\bfseries Candidato:}\\[0.2cm]
            Simone Garau
        \end{flushright}
    \end{minipage}

    \vfill
    {\large Anno Accademico 2024/2025 \par}
\end{titlepage}

\clearpage

\chapter*{Abstract}
% abstract.tex - Abstract della tesi
\chapter{Abstract}

L'ottimizzazione del \textit{Length of Stay} (LOS) ospedaliero è una sfida fondamentale per conciliare efficienza economica e qualità delle cure. Questa tesi, sviluppata all'interno del progetto TEXLOS \cite{texlos2026}, propone un sistema innovativo d'intelligenza artificiale per la classificazione retrospettiva e l'individuazione dei colli di bottiglia logistici responsabili di degenze critiche ($\ge 20$ giorni) partendo dai \textit{Log degli Eventi} clinici.

L'approccio metodologico supera i limiti del \textit{Process Mining} tradizionale introducendo il paradigma dello \textbf{Storytelling}: una traduzione semantica dei pattern cronologici in narrazioni in linguaggio naturale. Tali testi alimentano un classificatore Transformer (\texttt{bert-base-uncased}), addestrato mediante \textit{Focal Loss} per mitigare il severo sbilanciamento dei dati ospedalieri.

Avendo accertato l'apprendimento strutturale del modello, viene risolto il problema clinico della scatola nera ("black box"). Si applica un innovativo meccanismo \textsl{adattivo} di \textbf{Integrated Gradients Adattivi} -- fuso alla ricomposizione semantica intelligente dei token frammentati da BERT -- per interrogare a posteriori le scelte della rete. Assegnando inequivocabili "punteggi d'impatto" univoci direttamente alle specifiche azioni nosocomiali storiche, si fornisce ai decisori medici un cruscotto diagnostico di processo (heatmap) matematicamente rigoroso per risolvere empiricamente le inefficienze ospedaliere.


\tableofcontents

% --- Corpo della Tesi ---
\mainmatter
\onehalfspacing
\sloppy

% capitoli/01_introduzione.tex - Capitolo 1: Introduzione

\chapter{Introduzione}
\label{ch:introduzione}

\section{Contesto e Motivazione}
L'ottimizzazione dei processi sanitari rappresenta oggi una delle sfide principali per le amministrazioni ospedaliere, spinte dalla necessità di coniugare la qualità delle cure con la sostenibilità economica delle strutture. In questo contesto, il controllo e la previsione del \textit{Length of Stay} (LOS), ovvero la durata della degenza ospedaliera dei pazienti, rivestono un ruolo cruciale. Un LOS eccessivamente prolungato non solo grava sulle risorse finanziarie e strutturali dell'ospedale, come la disponibilità di posti letto e il carico di lavoro del personale, ma espone anche i pazienti a un maggiore rischio di infezioni nosocomiali e complicazioni cliniche.

Storicamente, l'analisi dei percorsi clinici si è avvalsa del \textit{Process Mining}, una disciplina che consente di estrarre conoscenza preziosa dai log degli eventi registrati dai sistemi informativi sanitari, esplorando e modellando i flussi di cura. Tuttavia, le metodologie classiche di Process Mining mostrano limiti nel gestire la complessità semantica e la sequenzialità non lineare intrinseche alle tracce cliniche dei pazienti, spesso frammentate e altamente variabili.

Negli ultimi anni, l'evoluzione dell'Intelligenza Artificiale ha aperto nuove prospettive, in particolare grazie all'introduzione dei modelli basati sull'architettura Transformer e all'avvento dei Large Language Models (LLM). L'intuizione fondamentale risiede nell'assimilare una sequenza di eventi clinici a una narrazione, traducendo le tracce in veri e propri documenti di testo (\textit{Storytelling}). I modelli linguistici avanzati, come BERT, possiedono infatti la capacità intrinseca di comprendere il contesto bidirezionale di una sequenza temporale, cogliendo dipendenze a lungo termine e sfumature che i modelli statistici tradizionali ignorano. 

Nonostante le elevate prestazioni predittive raggiunte dalle reti neurali profonde, la loro natura di ``scatola nera'' (\textit{black box}) costituisce un ostacolo significativo all'adozione pratica in ambito clinico, dove la fiducia degli operatori sanitari e la trasparenza decisionale sono requisiti imprescindibili. Al medico o all'amministratore ospedaliero non basta sapere \textit{se} un paziente avrà una degenza critica, ma serve comprendere il \textit{perché}, in modo da poter intervenire attivamente sulle inefficienze strutturali o sui colli di bottiglia dei processi. È qui che si inserisce l'Explainable AI (XAI), che si prepone di decostruire le decisioni dell'Intelligenza Artificiale fornendo motivazioni interpretabili per l'uomo.

\section{Obiettivi della Tesi}
Questa tesi si colloca all'interno del progetto di ricerca TEXLOS (\textit{Investigating the Impact of Outpatient Services on Length of Stay: an Easily Interpretable Approach}) e si pone due obiettivi primari fusi in un'unica pipeline di analisi.

Il primo obiettivo è lo sviluppo di un sistema di classificazione binaria basato sull'architettura Transformer per imparare a prevedere se la degenza di un paziente supererà una soglia critica, fissata a 20 giorni. Il modello fa leva sull'approccio dello \textit{Storytelling}, trasformando i log in formato XES in narrazioni testuali ed impiegando \texttt{bert-base-uncased} per catturare le connessioni logiche tra i vari eventi clinici.

Il secondo e più innovativo obiettivo riguarda l'interpretabilità del modello predittivo. Attraverso l'uso di tecniche avanzate di Explainable AI, in particolare l'algoritmo degli Integrated Gradients (IG) e le mappe di attenzione, il lavoro mira a estrapolare e quantificare l'impatto di ogni singola azione clinica sulla durata prolungata della degenza. L'obiettivo ultimo è fornire alle amministrazioni ospedaliere uno strumento diagnostico di processo capace di evidenziare in modo chiaro ed inequivocabile i colli di bottiglia, traducendo token sparsi in azioni cliniche leggibili provviste di un punteggio d'impatto rigoroso.

\section{Struttura dell'Elaborato}
Il presente documento è organizzato come segue:
\begin{itemize}
    \item Il \textbf{Capitolo \ref{ch:stato_arte}} esplora lo stato dell'arte, partendo dall'uso del Process Mining in ambito clinico fino all'impiego del Deep Learning e degli LLM per la classificazione di tracce. Particolare attenzione viene posta ai concetti di Explainable AI e all'approccio Storytelling, citando progetti pionieristici come LEGOLAS.
    \item Il \textbf{Capitolo \ref{ch:metodologia}} descrive in dettaglio i passaggi tecnici del lavoro: dall'estrazione e preprocessing dei dati alla creazione degli embedding, esponendo la configurazione della rete neurale (BERT), la gestione dello sbilanciamento estremo delle classi mediante Focal Loss e il setup sperimentale. Verrà anche discusso un esperimento esplorativo di data augmentation tramite LLM \textit{generativi}.
    \item Il \textbf{Capitolo \ref{ch:risultati}} riporta i risultati della classificazione in termini di accuratezza e metriche di valutazione. Successivamente, dettaglia criticamente l'implementazione adattativa degli Integrated Gradients e illustra il meccanismo di mappatura dei sub-token, mostrando concretamente come i risultati dell'algoritmo permettono di individuare i colli di bottiglia nel processo clinico.
    \item Il \textbf{Capitolo \ref{ch:conclusioni}} riassume i traguardi raggiunti dalla tesi, evidenzia i limiti metodologici intrinseci dell'approccio proposto, chiarendo la differenza tra correlazione e causalità, e traccia possibili traiettorie per gli sviluppi futuri.
\end{itemize}

% % capitoli/02_stato_arte.tex - Capitolo 2: Stato dell'Arte

\chapter{Stato dell'Arte}
\label{ch:stato_arte}

\section{Process Mining in Ambito Clinico}
Il settore sanitario genera quotidianamente moli di dati strutturati e non strutturati che riflettono l'interazione dei pazienti con le strutture ospedaliere, rendendo i \textit{Log degli Eventi} una miniera di informazioni fondamentale per analizzare oggettivamente i percorsi di cura. La gestione e il miglioramento di questi percorsi sono sfide affrontate storicamente attraverso il \textit{Process Mining}.

In ambito clinico, il Process Mining consiste in una serie di tecniche volte a estrarre modelli di processo dai dati (ad esempio tracce esportate in formato XES) per scoprire, monitorare e ottimizzare l'erogazione delle cure mediche e l'efficienza dei servizi. Gli algoritmi di Process Discovery o di Conformance Checking permettono alle amministrazioni di individuare anomalie procedurali e deviazioni dai percorsi terapeutici codificati (Clinical Pathways). Sebbene questi metodi godano di robustezza matematica per mappare il sistema, si sono spesso rivelati rigidi di fronte alle variabili caotiche e all'alta stocasticità dei flussi ospedalieri, che esibiscono dipendenze temporali complesse difficili da cogliere attraverso soli grafi deterministici.

\section{Evoluzione del Deep Learning: Da RNN a Transformer}
La necessità di mitigare le carenze del Process Mining basato su regole si è tradotta nell'integrazione di tecniche di Data Mining esplorativo, giungendo ben presto all'utilizzo del \textit{Deep Learning}. 

Inizialmente, l'esigenza di elaborare tracce cliniche di lunghezza variabile ha portato all'impiego delle architetture ricorrenti, in primis Recurrent Neural Networks (RNN) e Long Short-Term Memory (LSTM), poiché naturalmente predisposte per dati sequenziali. Tuttavia, il difetto cronico delle RNN legato alla dispersione del gradiente (\textit{vanishing gradient}) e l'incapacità intrinseca di gestire parallelismi e interdipendenze a lungo raggio in storie cliniche complesse ne hanno fortemente limitato l'efficacia pratica.

Il cambio di paradigma epocale avviene nel 2017 con l'introduzione dell'architettura Transformer e l'invenzione dei meccanismi di \textit{Self-Attention}. Un passo decisivo verso l'analisi semantica delle tracce ospedaliere è stato il superamento delle classiche codifiche One-Hot, tipiche dell'era pre-Transformer, in favore di \textit{embeddings} linguistici. L'avvento dei Large Language Models (LLM), tra cui Bidirectional Encoder Representations from Transformers (BERT) \cite{devlin2019bert}, ha consentito di addestrare vettori profondamente contestuali su moli sterminate di dati testuali non etichettati, estraendo conoscenze sintattiche e logiche che vengono trasferite con successo a task applicativi specifici. Nel caso clinico, ciò si traduce in un modello capace di osservare e valutare un'intera traccia di eventi contemporaneamente e in modo bidirezionale, anziché passo dopo passo nel tempo.

\section{L'Approccio Storytelling: Il Progetto LEGOLAS}
Affinché i modelli Transformer, ed in particolare BERT, possano elaborare con profitto i dati tratti da log storici di Process Mining, si è resa necessaria una forte innovazione nella tecnica di rappresentazione dell'informazione (\textit{embedding}). Nasce da questa esigenza l'approccio \textit{Storytelling}.

L'idea fondamentale, che ispira anche il presente lavoro di tesi, si radica nel recente e pionieristico progetto \textit{LEGOLAS} (Leveraging a Large Language Model to Predict Hospital Admissions of Emergency Department Patients) \cite{PASQUADIBISCEGLIE2025128224}. L'intuizione principale di LEGOLAS consiste nel superare la mera concatenazione meccanica di eventi e timestamp discreti per spingersi nella generazione di veri e propri documenti di testo coerenti, in linguaggio naturale. 

La traccia di un paziente non è più vissuta dall'intelligenza artificiale come una sequela di array numerici estratti dal registro dell'ospedale, ma viene trascritta in sotto forma di "storia" cronologica, in modo che il testo contenga esplicitamente la semantica degli eventi temporali e delle distanze tra un'azione clinica e l'altra. Poiché BERT è un LLM pre-addestrato sul linguaggio naturale umano, nutrirlo con "racconti clinici" formattati linguisticamente ad hoc permette un salto prestazionale considerevole in fase di Fine Tuning per task di classificazione, mitigando problemi di dimensionalità esplosiva tipici del Process Mining.

\section{Explainable AI (XAI) per i Modelli Transformer}
Per superare i timori sulla natura a scatola nera delle reti neurali, l'Explainable AI (XAI) emerge come necessità non eludibile. Nel contesto del Deep Learning per il settore medico, gli algoritmi trasparenti come il Decision Tree non reggono in termini di performance rispetto al Deep Learning, rendendo cruciale sviluppare metodi post-hoc (che si applicano a modello addestrato) per interpretare reti complesse.

Nel caso dei modelli basati su Transformer, si analizza il focus cognitivo della rete ricorrendo comunemente alle \textit{Mappe di Attenzione} della rete stessa, un proxy del processo di ragionamento interno del modello. Tuttavia, le \textit{Attention Maps} riflettono i "pesi" posizionali generici, che non indicano un contributo deduttivo diretto verso una specifica decisione di classificazione, e sono note per essere affette da \textit{noise} e scarsa proporzionalità reale all'importanza di un token.

Al fine di ottenere metriche di feature importance rigorose rispetto alla classe predetta (es. "degente lungo termine"), la letteratura scientifica più recente suggerisce di guardare al calcolo dei Gradienti Integrati. Proposto originariamente da Sundararajan et al. \cite{sundararajan2017axiomatic}, il metodo degli \textit{Integrated Gradients} (IG) è una tecnica assiomatica e robusta. L'intuizione dell'algoritmo IG consiste nell'accumulare i gradienti lungo una linea retta nello spazio delle feature, passando gradualmente da un input "baseline" neutro fino all'input reale testuale (la storia del paziente), garantendo matematicamente che la somma del contributo di tutti i singoli token corrisponda esattamente allo scarto tra la probabilità emessa sulla baseline neutra e quella finale calcolata. Questa aderenza matematica (\textit{Completeness}) fa degli IG lo strumento d'elezione per misurare minuziosamente quale parola ha spinto o inibito la decisione del classificatore BERT.

% % capitoli/03_metodologia.tex - Capitolo 3: Metodologia

\chapter{Metodologia}
\label{ch:metodologia}

\section{Estrazione e Preprocessing dei Dati}
Il punto di partenza dell'esperimento è costituito da log ospedalieri estratti in formato XES. L'obiettivo primario di questa fase iniziale è strutturare e formattare tali log non strutturati per renderli digeribili all'architettura neurale sottostante il Transformer. 

\subsection{Analisi Esplorativa del Dataset (EDA)}
Per acquisire contezza della natura e del volume dei log estratti pre-elaborazione, è stata condotta un'Analisi Esplorativa dei Dati (EDA) sviluppando una pipeline automatizzata in Python (basata principalmente sulla libreria \texttt{pandas} per l'efficiente gestione e \textit{chunking} massivo delle estrazioni). Il dataset grezzo comprende 7.393 casi clinici unici (singoli pazienti ospitati), all'interno dei quali sono registrati 88.115 eventi granulari totali non etichettati.

Si rileva contestualmente l'ingente sbilanciamento delle classi, cardine della sfida analitica affrontata in questo progetto: su 7.393 pazienti totali, ben 6.584 esibiscono una degenza classificabile come normale (classe 0), e solamente 809 costituiscono campioni d'anomalia con degenza critica $\ge 20$ giorni (classe 1).

\begin{figure}[htbp]
    \centering
    \includegraphics[width=0.8\textwidth]{"img/actions_per_case_distribution.png"}
    \caption{Distribuzione statistica globale del numero di azioni processuali registrate per singolo caso clinico. Il grafico si caratterizza per una spiccata asimmetria (\textit{right-skewed}), evidenziando come una concentrazione primaria di interazioni ospedaliere rapide all'ingresso si contrapponga ad una ``lunga coda'' indicante pazienti complessi trattenuti in cure iterative stratificate.}
    \label{fig:actions_distribution}
\end{figure}
\FloatBarrier

Come illustrato in Figura \ref{fig:actions_distribution}, il numero di azioni per paziente presenta una forte asimmetria destra (\textit{right-skewness}): la media è di 12.92 azioni, mentre la mediana scende a 9.00, segnalando che la distribuzione è tirata verso l'alto da una ``lunga coda'' (\textit{long tail}) di pazienti complessi con decine o centinaia di re-interazioni.

\begin{figure}[htbp]
    \centering
    \includegraphics[width=0.8\textwidth]{"img/class_quantity_boxplot.png"}
    \caption{Boxplot di confronto diretto tra il numero d'interazioni erogate ai pazienti di classe 0 (degenza limitata) e di classe 1 (degenza prolungata). Risalta ad occhio nudo la traslazione dei quartili a sfavore della classe anomala, gravata intrinsecamente da percorsi ben più caotici, iterativi ed onerosi quantitativamente per la clinica.}
    \label{fig:class_quantity_boxplot}
\end{figure}
\FloatBarrier

Tale diversità emerge chiaramente analizzando le durate per classe (Figura \ref{fig:class_quantity_boxplot}): appare evidente il divario numerico, con la classe 1 (long-LOS) caratterizzata da percorsi mediani molto più lunghi e stratificati rispetto ai pazienti con degenti brevi. Questa evidenza giustifica l'adozione di un'architettura avanzata come BERT. A differenza di una stringa Markoviana classica che frammenterebbe i calcoli su sequenze enormi, l'architettura basata su \textit{Self-Attention} garantisce un'elaborazione in parallelo per mappare le dipendenze logiche tra decine di step ospedalieri.

\begin{figure}[htbp]
    \centering
    \includegraphics[width=\textwidth]{"img/class_action_prevalence_diff_top.png"}
    \caption{Panoramica estrapolata dei \textit{delta} di prevalenza percentuale (sbilanciamento) per le singole azioni ospedaliere fortemente correlate in modo endogeno alla classe d'emergenza logistica 1.}
    \label{fig:class_action_prevalence}
\end{figure}
\FloatBarrier

Un'ulteriore esplorazione evidenzia la distribuzione asimmetrica delle frequenze delle singole procedure (Figura \ref{fig:class_action_prevalence}). Risultano particolarmente pronunciati gap associati ad alta probabilità critica per eventi clinici come la ``Visita Infettivologica'', la ``Radiografia Torace a Letto'' o il ``Posizionamento Catetere Venoso Centrale''. Pur trattandosi di statistiche esplorative prive di valore causale formale, tali osservazioni anticipano i pattern che il modello imparerà a pesare durante il training e che verranno interpretati tramite XAI nel Capitolo \ref{ch:risultati}.

\subsection{Approcci e Formattazione degli Embedding}
Poiché il task ingegneristico mira strettamente ad accoppiare eventi clinici elaborati a corretto formato ed interfacciandoli poi col \textit{Learning} del Transformer, è stata condotta una meticolosa opera preliminare testando tre filosofie distinte di \textit{embedding}:

\begin{enumerate}
    \item \textbf{Azioni cliniche + Binning Temporale:} In questo approccio, il preprocessing concatena coppie formate dall'azione clinica e da un token temporale speciale. Il timing continuo tra le azioni viene discretizzato (\textit{binning}) in categorie predefinite. Un esempio di traccia processata è il seguente:
\begin{minted}[fontsize=\small]{json}
{"case_id": "5870415", 
 "tokens": ["Complete abdomen ultrasound", "_TIME_DELTA_START_", 
            "Computerized tomography of the skull", 
            "_TIME_DELTA_START_", "Complete abdomen ultrasound", 
            "_TIME_DELTA_LONG_", "Computerized tomography of the skull", 
            "_TIME_DELTA_MEDIUM_"], 
 "label_id": 0
}
\end{minted}

    La scelta delle soglie di discretizzazione non è arbitraria ma \textit{data-driven}: le categorie sono state definite analizzando la distribuzione empirica dei delta temporali nel dataset, riportata in Tabella~\ref{tab:time_delta_distribution}. La distribuzione rivela una struttura bimodale con la grande maggioranza degli intervalli concentrata entro la finestra dei 7 giorni (97.0\%), e una coda esigua di eventi molto dilazionati. Le soglie sono state quindi fissate sui breakpoint clinici naturali --- un'ora, un giorno, sette giorni e trenta giorni --- che separano fasi di cura operativamente distinte: accertamenti rapidi, degenza standard, ricovero prolungato e cronicità.

    \begin{table}[htbp]
        \centering
        \resizebox{\textwidth}{!}{%
        \begin{tabular}{|l|c|r|r|}
            \hline
            \textbf{Token} & \textbf{Intervallo} & \textbf{Conteggio} & \textbf{Freq.} \\
            \hline
            \texttt{\_TIME\_DELTA\_START\_}      & $\Delta = 0$ s                         & 5\,842  &  6.9\% \\
            \texttt{\_TIME\_DELTA\_IMMEDIATE\_}  & $0 < \Delta \le 1\text{h}$             & 19\,456 & 23.1\% \\
            \texttt{\_TIME\_DELTA\_SHORT\_}      & $1\text{h} < \Delta \le 1\text{gg}$   & 29\,887 & 35.4\% \\
            \texttt{\_TIME\_DELTA\_MEDIUM\_}     & $1\text{gg} < \Delta \le 7\text{gg}$  & 26\,649 & 31.6\% \\
            \texttt{\_TIME\_DELTA\_LONG\_}       & $7\text{gg} < \Delta \le 30\text{gg}$ &  2\,296 &  2.7\% \\
            \texttt{\_TIME\_DELTA\_VERY\_LONG\_} & $\Delta > 30\text{gg}$                 &    233  &  0.3\% \\
            \hline
            \multicolumn{2}{|l|}{\textbf{Totale delta calcolati}} & 84\,363 & 100\% \\
            \hline
        \end{tabular}%
        }
        \caption{Distribuzione empirica dei delta temporali nel dataset e corrispondente token di discretizzazione assegnato. Il totale è inferiore agli 88.115 eventi grezzi poiché il primo evento di ciascuna traccia non possiede un evento precedente da cui calcolare il delta.}
        \label{tab:time_delta_distribution}
    \end{table}
    \FloatBarrier

    La funzione \texttt{categorize\_time\_delta\_vectorized}, implementata con \texttt{numpy} per un processing completamente vettorializzato sull'intero dataset, traduce i delta continui nei token speciali secondo questa logica:

\begin{minted}[fontsize=\small]{python}
def categorize_time_delta_vectorized(deltas):
    conditions = [
        deltas == 0,                               # Delta nullo: eventi simultanei
        (deltas > 0) & (deltas <= 3600),           # 0-1 ora:   Immediato
        (deltas > 3600) & (deltas <= 86400),       # 1h-1gg:    Breve
        (deltas > 86400) & (deltas <= 604800),     # 1-7gg:     Medio
        (deltas > 604800) & (deltas <= 2592000),   # 7-30gg:    Lungo
    ]
    choices = [
        '_TIME_DELTA_START_',       # delta = 0
        '_TIME_DELTA_IMMEDIATE_',   # 0 < delta <= 1h
        '_TIME_DELTA_SHORT_',       # 1h < delta <= 1gg
        '_TIME_DELTA_MEDIUM_',      # 1gg < delta <= 7gg
        '_TIME_DELTA_LONG_',        # 7gg < delta <= 30gg
    ]
    return np.select(conditions, choices, default='_TIME_DELTA_VERY_LONG_')
\end{minted}

    \item \textbf{Array concatenati (Azioni + Distanze numeriche):} Questa seconda forma di preprocessing costruisce l'embedding affiancando due array distinti: uno contenente esclusivamente la sequenza testuale delle azioni cliniche e un array parallelo numerico (\textit{float}) esprimente la distanza in secondi dall'evento precedente. Un esempio:
\begin{minted}[fontsize=\small]{json}
{"case_id": "5870415", 
 "tokens": ["Complete abdomen ultrasound", "Computerized tomography of the skull", 
            "Complete abdomen ultrasound", "Computerized tomography of the skull"], 
 "time_deltas": [0.0, 0.0, 728160.0, 178980.0], 
 "label_id": 0
}
\end{minted}

    \item \textbf{Storytelling in linguaggio naturale:} L'approccio finale, che costituisce il core metodologico del progetto TEXLOS \cite{texlos2026} e che è stato preferito ai precedenti, si rifà al pionieristico lavoro di traduzione semantica dei log del progetto LEGOLAS \cite{PASQUADIBISCEGLIE2025128224}. Anziché alimentare il classificatore con una combinazione mista di vettori numerici e codici, gli eventi e i vincoli cronologici logici della traccia sono tradotti letteralmente in frasi inglesi strutturate grammaticalmente, producendo interi paragrafi coerenti narranti il percorso diagnostico/terapeutico di ciascun paziente.
\begin{minted}[fontsize=\small]{json}
{"case_id": "5870415", 
 "story": "The patient, 87 years old, underwent a series of examinations and 
 interventions during hospitalization.
After 0 seconds, the following examinations were performed simultaneously: 
Complete abdomen ultrasound and Computerized tomography of the skull.
After 728,160 seconds, Complete abdomen ultrasound was performed.
After 178,980 seconds, Computerized tomography of the skull was performed.",
"label_id": 0
}
\end{minted}
\end{enumerate}

\subsection{Dettagli Implementativi dell'Algoritmo di Storytelling}
Per colmare la distanza che intercorre tra i log estratti in formato grezzo e i testi discorsivi forniti in addestramento, è stata sviluppata un'infrastruttura di \textit{parsing} ed iniezione in linguaggio Python. Dal punto di vista informatico, il processo si divide in due fasi computazionali: la strutturazione in memoria del file XES e la generazione cronologica della storia narrata. 

Nella prima fase entra in gioco il modulo \texttt{xes\_parser.py}. Affrontando l'ingente volume di eventi clinici asincroni tipici di un'analisi di Process Mining, l'algoritmo si destreggia tra la libreria per l'ingestione dei log (\texttt{pm4py}) e le più ottimali e leggere architetture dati di \texttt{pandas}. Quest'ultimo processa massivamente il dataset suddividendo mediante \textit{groupby} e parallelismi nativi le righe afferenti ai medesimi pazienti; si procede poi convertendo in formato standard i \textit{timestamp} (le date d'esecuzione dell'atto medico), le stringhe degli attori coinvolti ed instradando così ogni frammento cronologico entro uno strato informativo solido (oggetto \textit{PatientTrace}).  

La fase autoriale vera e propria è governata dalla classe \texttt{StoryGenerator} (all'interno dello script \texttt{story\_generator.py}). Le specifiche di traduzione che intercettano queste \textit{trace} vertono su tre stringenti requisiti ingegneristici:
\begin{itemize}
    \item \textbf{Gestione del Tempo (Time Deltas):} L'evoluzione clinica viene misurata estraendo rigorosamente il \textit{delta temporale in secondi} che separa l'azione attuale dalla precedente operazione registrata. Questo gap continuo viene iniettato nel template e calcolato esplicitamente mediante un blando \textit{parsing} aritmetico (\texttt{current\_time - prev\_time}), oggettivando la scansione temporale. Nessuna approssimazione logica interviene nel testo finale; i secondi marcano senza preconcetti le tempistiche di risposta dei triage o la dilatazione cronica della degenza (si passa da "0 secondi" per l'esordio iniziale, a svariate centinaia di migliaia di secondi per tac e risonanze d'esito fatte la settimana d'uscita).
    \item \textbf{Gestione degli Eventi Simultanei:} Non è inusuale, nei processi su scala giornaliera densi come in ospedale, rilevare molteplici interazioni registrate col medesimo identico \textit{timestamp}. Un esempio classico si verifica con batterie di check-up del sangue incodificate alla stessa ora informatica dal personale. Per evitare grave ridondanza lessicale che sovraccaricherebbe indebitamente l'asse d'attenzione di BERT con frasi ripetute, il generatore intercetta internamente gli aggregati (tramite la funzione \texttt{\_group\_simultaneous\_events}), fondendoli all'interno di specifici \textit{bullet points} o paragrafi discorsivi compatti nel seguente formato: \textit{"After $X$ seconds, the following examinations were performed simultaneously: Action A, and Action B"}.
    \item \textbf{Indipendenza dal Protocollo di Triage:} Il generatore proposto si differenzia strutturalmente dal progetto LEGOLAS, il quale era progettato per interpolare un insieme fisso di variabili fisiologiche standard da Pronto Soccorso (frequenza respiratoria, temperatura corporea, scala del dolore, ecc.), compilando un template narrativo predefinito. Nel contesto di TEXLOS, i pazienti attraversano reparti eterogenei con procedure cliniche non standardizzate, prive di una nomenclatura universale codificata. Per gestire questa variabilità, il generatore adotta dizionari interni flessibili e \textbf{regole di template matching dinamiche}, agnostiche rispetto alla tipologia di azione clinica, che mappano la tassonomia locale dell'ospedale ai costrutti discorsivi del Transformer senza vincolare l'algoritmo a un sottoinsieme predefinito di sensori o misurazioni.
\end{itemize}

\newpage
A scopo esemplificativo, il seguente frammento di pseudo-codice descrive la routine adoperata dal software in \texttt{story\_generator.py} per snellire testualmente (\textit{trimming} sintattico) il sovraffolamento degli stadi simultanei ed iniettare i \textit{time-delta} per una traccia degenziale prima dell'invio in NLP ad estrapolazione avvenuta:

\begin{minted}[fontsize=\small]{python}
def create_narrative_from_trace(trace_events):
    sorted_events = sort_by_timestamp(trace_events)
    grouped_events = group_simultaneous_events(sorted_events)
    
    narrative_paragraphs = []
    prev_time = grouped_events[0].timestamp
    
    for group in grouped_events:
        time_elapsed = calculate_seconds_diff(prev_time, group.timestamp)
        
        if len(group) == 1:
            # Singolo evento testuale formattato
            activity_desc = apply_medical_template(group[0].activity)
            text = f"After {time_elapsed} seconds, {activity_desc} was performed."
        else:
            # Piu' azioni collassate per alleggerimento dell'Attention
            activities = [apply_medical_template(e.activity) for e in group]
            concat_acts = join_with_commas_and(activities)
            text = f"After {time_elapsed} seconds, the following examinations " \
                   f"were performed simultaneously: {concat_acts}."
                   
        narrative_paragraphs.append(text)
        prev_time = group.timestamp # Sposta l'ancora
        
    return "\n\n".join(narrative_paragraphs)

\end{minted}

\section{Architettura del Modello e Limitazioni dei BERT Clinici}
Alla base del sistema di previsione giace \texttt{bert-base-uncased}, un Large Language Model reso disponibile pubblicamente sulla piattaforma HuggingFace. Questo Transformer, limitato alla componente Encoder, viene sottoposto a Fine Tuning aggiungendo al termine dell'architettura standard una testa di classificazione (\textit{classification head}) volta a predire la natura binaria del problema: degenza normale (sotto i 20 giorni) contro degenza critica (uguale o superiore a 20 giorni).

Il primo istinto sperimentale è stato chiaramente quello di testare modelli \textit{Domain-Specific}, confrontando le prestazioni di varianti pre-addestrate unicamente su corpus clinici -- quali \texttt{biobert}, \texttt{pubmedBert}, \texttt{blue-bert} e \texttt{clinical-bert} -- con quelle di \texttt{bert-base-uncased}. L'ipotesi sottostante era che un vocabolario medico radicato \textit{ab initio} avrebbe accelerato la decodifica delle diagnosi. L'evidenza empirica ha sorprendentemente sancito prestazioni sovrapponibili tra i vari modelli e, anzi, una \textit{Balanced Accuracy} leggermente inferiore per i modelli clinici rispetto al canonico modello generalista. 

Tale anomalia si spiega analizzando l'altissima specificità terminologica del dataset della tesi: i modelli addestrati su dataset clinici (come MIMIC) sono spesso altamente specifici e ottimizzati per una formale nomenclatura standard inglese. Nel nostro caso, il nomenclatore locale dell'ospedale era ricco di acronimi, contrazioni proprietarie e formulazioni non standardizzate, vanificando almeno parzialmente il vantaggio della pregressa conoscenza medica dei pesi. Questo disallineamento è critico soprattutto a livello di \textit{tokenizer}: i tokenizer dei vari modelli clinici, agendo su un vocabolario molto più verticale, frammentavano le parole spurie del dataset della tesi perdendo efficacia. Di contro, il tokenizer base di \texttt{bert-base-uncased} si è dimostrato molto più elastico, resiliente e capace di preservare sufficiente compattezza semantica pur al cospetto di un dominio estraneo.

\section{Giustificazione Architetturale: Analisi del Trade-off Computazionale}
In fase di definizione dell'architettura neurale, è stata condotta un'indagine esplorativa scalando il modello base verso la sua variante maggiorata \texttt{bert-large-uncased} (circa 340 milioni di parametri contro i 110 milioni originari). L'intento era tracciare la scalabilità dell'accuratezza su una rete più stratificata.

Un test preliminare su una singola esecuzione ha restituito metriche di rilievo sulla classe critica (LOS $\ge 20$ giorni), registrando una Precision del 73.65\% e una Recall dell'84.36\%. Tuttavia, pur evidenziando un lieve incremento prestazionale grezzo rispetto alla versione base, l'esperimento ha fatto emergere severi limiti di scalabilità e di \textit{data plateauing}. Da un lato, l'aumento dimensionale del modello non ha prodotto un salto di qualità proporzionale al costo computazionale, suggerendo che un bacino di circa 7.000 tracce cliniche non sia sufficiente a saturare la capacità di astrazione di un'architettura "Large", portando di fatto la rete in una fase di plateau informativo.

Dall'altro lato, i vincoli hardware hanno reso la sperimentazione ingegneristicamente insostenibile. La saturazione della VRAM ha imposto una drastica riduzione della \textit{batch size}, precipitata dai 22 campioni gestibili dal modello base ad appena 4 campioni per il modello large. Questo crollo ha dilatato esponenzialmente i tempi di addestramento, rendendo materialmente inapplicabile la \textit{Stratified 5-Fold Cross Validation}, requisito imprescindibile delineato in questa tesi per garantire la robustezza statistica del classificatore dinanzi a uno sbilanciamento di classi così severo.

Di conseguenza, il modello \texttt{bert-base-uncased} è stato eletto come "ottimo paretiano" del progetto, offrendo il compromesso ideale tra capacità predittiva, sostenibilità hardware e garanzia di validazione metodologica rigorosa.

\section{Focal Loss e Sbilanciamento delle Classi}
Il dataset presenta uno dei paradossi strutturali più frequenti della Process Analytics in sanità: lo \textit{sbilanciamento estremo delle classi}. Di fatto, le casistiche in cui il paziente sfiora o sfora una degenza continua di venti giorni (classe 1 - Critico) costituiscono una netta e fisiologica minoranza anomala nei log degli eventi rispetto alla travolgente maggioranza dei passaggi di routine o dimissioni rapide (classe 0 - Normale).

Una prima fase di addestramento gestita dalla standard \textit{Binary Cross Entropy} (BCE) ha confermato l'incapacità del modello di apprendere in modo bilanciato: il classificatore tendeva a privilegiare la classe dominante, producendo di fatto una predizione prossima al \textit{ZeroR} e penalizzando la sensibilità sui casi critici — esattamente quelli di primario interesse per le amministrazioni ospedaliere.

Per correggere questo bias strutturale e ricalibrare il contributo al gradiente in modo differenziale, è stata adottata in fase di fine-tuning la funzione di costo denominata \textbf{Focal Loss} \cite{lin2017focal}. Questa funzione mutua la sua origine nel dominio della \textit{Computer Vision}, dove nacque per correggere il massiccio sbilanciamento tra uno sconfinato spazio di \textit{background} e piccoli oggetti da identificare. Applicata al processing di tracce cliniche, la Focal Loss riduce il peso computazionale (\textit{loss contribution}) degli esempi facili e sovrabbondanti, in favore dei campioni mal classificati.

\subsection{Formulazione Matematica della Focal Loss}
Il punto di partenza è la funzione di costo \textit{Cross Entropy} (CE) per la classificazione binaria. Definito $y \in \{\pm 1\}$ come la classe ground-truth e $p \in [0, 1]$ come la probabilità stimata dal modello per la classe $y=1$, la formulazione classica è:

\begin{equation}
CE(p, y) = \begin{cases} -\log(p) & \text{se } y=1 \\ -\log(1-p) & \text{altrimenti.} \end{cases}
\end{equation}

Per convenienza notazionale, definendo $p_t$:
\begin{equation}
p_t = \begin{cases} p & \text{se } y=1 \\ 1-p & \text{altrimenti,} \end{cases}
\end{equation}
la CE si può riscrivere compattamente come $CE(p_t) = -\log(p_t)$.

Un metodo comune per mitigare lo sbilanciamento puro introduce un fattore di peso $\alpha \in [0, 1]$ per la classe 1 e $1-\alpha$ per la classe -1. Definendo $\alpha_t$ in modo analogo a $p_t$, si ottiene la \textit{$\alpha$-balanced Cross Entropy}:
\begin{equation}
CE(p_t) = -\alpha_t \log(p_t)
\end{equation}

Tuttavia, pur bilanciando l'importanza numerica tra positivi e negativi, questa formulazione non differenzia tra esempi \textit{facili} (\textit{easy negatives}) ed esempi \textit{difficili} (\textit{hard examples}). Per abbattere il peso sproporzionato degli innumerevoli esempi facili (i pazienti con degenze standard facilmente prevedibili), Lin et al. introducono un fattore modulante $(1-p_t)^\gamma$, dove $\gamma \ge 0$ è un iperparametro noto come \textit{parametro di focalizzazione} (\textit{focusing parameter}). La formulazione completa della \textbf{Focal Loss} (nella sua variante $\alpha$-balanced pratica) diviene quindi:

\begin{equation}
FL(p_t) = -\alpha_t (1-p_t)^\gamma \log(p_t)
\end{equation}

Questa funzione gode di proprietà analitiche fondamentali per il nostro task:
\begin{itemize}
    \item Se un esempio ospedaliero è mal classificato o raro (quindi $p_t$ è piccolo), il fattore modulante $(1-p_t)^\gamma$ tende a 1 e la \textit{loss contribution} rimane inalterata, garantendo l'apprendimento.
    \item Se un esempio è facile e ben classificato ($p_t \rightarrow 1$), il fattore modulante tende a 0, penalizzando (\textit{down-weighting}) fortemente l'impatto sul gradiente globale.
    \item Il parametro $\gamma$ regola il tasso di questa penalizzazione (addolcendo l'iper-dominanza della classe 0 nel nostro dataset), esonerando l'analista dalla necessità di applicare euristiche forzate e distruttive a valle, come un aggressivo \textit{undersampling} che distruggerebbe preziosi pattern temporali globali.
\end{itemize}

L'effetto della parametrizzazione $\gamma$ si evince analizzando visivamente l'andamento della funzione per vari scenari, confrontando la Focal Loss con la Cross Entropy standard (Figura \ref{fig:focal_loss_comparison}).

\begin{figure}[htbp]
    \centering
    \includegraphics[width=0.8\textwidth]{"img/focal loss fig 1.png"}
    \caption{Confronto tra la Cross Entropy (CE) e diverse configurazioni della Focal Loss. Sull'asse delle ascisse è riportata la probabilità effettiva stimata dal modello, $p_t \in [0, 1]$, mentre sull'asse delle ordinate figura il valore della funzione di costo. Si noti come, al crescere dell'iperparametro di focalizzazione $\gamma$ (es. $\gamma=2, 5$), la \textit{loss} addebitata agli esempi facili e ben classificati ($p_t \rightarrow 1$) subisca un drastico schiacciamento rispetto al decadimento più lento della CE. Questo espediente permette all'osservazione di focalizzarsi dinamicamente sugli esempi più critici ed errati.}
    \label{fig:focal_loss_comparison}
\end{figure}
\FloatBarrier

Una volta garantita la capacità discriminatoria del modello su classi sbilanciate, si procede all'analisi della sua interpretabilità tramite le tecniche descritte nel capitolo successivo.

\section{Setup Sperimentale e Hyperparameter Tuning}
L'intera pipeline di esperimenti e addestramento supervisionato è stata materialmente elaborata e portata a convergenza su un setup hardware performante, incentrato su una workstation munita di CPU AMD Ryzen 9 5900X, 64 GB di RAM e GPU NVIDIA GeForce RTX 4090 (consentendo l'aumento efficiente delle batch size e l'addestramento distribuito del grafo tensoriale Transformer).

Per preservare l'integrità e il realismo della diagnostica probabilistica al cospetto dello sbilanciamento delle classi, la valutazione accademica si basa unicamente su iterazioni cicliche di \textit{Stratified K-Fold Cross Validation}. Questo garantisce che ogni singola split prodotta preservi identiche proporzioni di pazienti anomali/ordinari rispetto al database globale, smaltendo bias d'addestramento.

L'impostazione degli iperparametri del modello neurale è sfociata da un processo combinato: 
\begin{itemize}
    \item L'esplorazione automatica bayesiana condotta massivamente dal framework \texttt{Optuna}.
    \item Un minuzioso \textit{tuning manuale} a posteriori. Affinché i risultati non scivolassero in asimmetrie dimensionali non volute, le alterazioni manuali degli intervalli innescati da Optuna si sono preoccupate di preservare costanti le "proporzioni architetturali" auree del transformer, quali i rapporti intrinseci tra hidden layers, attention heads e scaling rate dello stack.
\end{itemize}

\subsection{Convergenza del Modello e Early Stopping}
A seguito del setup definito, l'addestramento della configurazione finale \texttt{bert-base-uncased} ha evidenziato un comportamento stabile. Come illustrato nella Figura \ref{fig:loss_curve}, l'andamento mostra una discesa costante della \textit{training loss}, mentre la \textit{validation loss} raggiunge il suo minimo all'epoca 10. Nelle epoche successive la curva di validazione inizia a divergere, indicando l'innesco dell'\textit{overfitting}. Grazie alla tecnica dell'\textit{early stopping} (impostata con un $\Delta$ minimo di $1\times 10^{-3}$ su una \textit{patience} di 5 epoche), il training si è interrotto all'epoca 15; i pesi del modello sono stati quindi ripristinati all'ottimo dell'epoca 10, garantendo la massima capacità di generalizzazione sui dati non visti per la successiva fase di test.

% STORIA LOSS
\begin{figure}[htbp]
    \centering
    \includegraphics[width=0.8\textwidth]{img/train_val_loss.png}
    \caption{Curva di apprendimento (Training vs Validation Loss) del modello definitivo durante il processo di \textit{fine-tuning}.}
    \label{fig:loss_curve}
\end{figure}
\FloatBarrier

\section{Indagine Esplorativa: Data Augmentation con LLM}
Nel tentativo di robustire l'astrazione di BERT sulla classe minoritaria (long-LOS), si è tentato un originale esperimento di estensione del dataset (\textit{Data Augmentation}) avvalendosi dell'intelligenza generativa degli attuali LLM generativi (decoder-only).

La pipeline esplorativa prevedeva l'impiego locale di modelli generativi allo stato dell'arte, nello specifico \texttt{phi4:14b} e \texttt{nous-hermes2:34b}, eseguiti tramite l'ambiente open-source \texttt{Ollama}. Per inibire le allucinazioni e circoscrivere l'effusione creativa dell'esposizione testuale, la \textit{Temperature} è stata fissata al minimo rigido ($T=0.0$). 

Il layer di codice Python istruiva severamente le reti a limitarsi ad un mero lavoro di riformattazione narrativa, passando iterativamente le coordinate base della traccia log:

\begin{minted}[fontsize=\scriptsize]{python}
prompt = f"""You are an expert physician who writes discursive clinical
narratives in English for BERT/NLP analysis.

IMPORTANT - MANDATORY MEDICAL TERMINOLOGY:
You must use EXACTLY the medical terminology specified in the clinical vocabulary.

PATIENT DATA:
- Information: {patient_desc}

MEDICAL EVENTS SEQUENCE (already standardized):
{events_text}

SPECIFIC RULES:
1. USE EXACTLY the provided activity names (already standardized)
2. Introduction with age and medical context
3. Separate paragraph for each temporal group
4. Format numbers: "13,380 seconds" (with commas for thousands)
5. Concurrent events: "the following examinations were performed simultaneously:"
6. Connectors: "The [activity] was performed", "After X seconds [activity] was executed"
7. Paragraphs separated by blank line
8. Keep medical terminology IDENTICAL to provided terminology
9. Always conclude with "At the end of the clinical pathway, the medical events were concluded."
10. NO translations or paraphrases of medical activity names
11. All text must be in ENGLISH"""
\end{minted}

Nonostante l'imbragatura direzionale e sintattica imposta a livello di framework, l'ingestione e la rielaborazione da parte delle reti ha portato a delle devianze gravose. A titolo di dimostrazione, viene qui incollata la narrativa sintetica prodotta dal modello \texttt{phi4:14b} a fronte del caso ospedaliero ID "5870415" (contenente originariamente 4 atti nosocomiali primari ed oggettivi estratti dai registri):

\begin{minted}[fontsize=\scriptsize]{json}
{
  "case_id": "5870415",
  "story_text": "The patient, 87 years old, underwent a series of examinations...
The Complete abdomen ultrasound was performed as the initial examination to 
assess abdominal structures for any abnormalities or pathologies that might 
require further investigation... This non-invasive imaging technique provided 
valuable insights... After 0 seconds, the following examinations were performed 
simultaneously: Computerized tomography of the skull and Complete abdomen ultrasound. The Coputerized tomography of 
the skull was conducted to evaluate any potential intracranial issues such as 
hemorrhages, tumors, or other abnormalities... After 728,160 seconds 
(approximately 8 days and 7 hours), another Complete abdomen ultrasound was 
executed. This follow-up examination aimed to monitor any changes...
At the end of the clinical pathway, the medical events were concluded."
}
\end{minted}

L'esempio documenta i limiti intrinseci nella generazione del testo impiegata per questo scopo predittivo. Nonostante le rigide indicazioni tecniche nel prompt testuale, il modello \texttt{phi4:14b} ha generato un'ampia \textbf{allucinazione semantica}, arricchendo indebitamente l'output testuale: il LLM ha aggiunto descrizioni procedurali e teoriche arbitrarie sul motivo degli esami (\textit{"to evaluate any potential intracranial issues such as hemorrhages"}), che risultano rumore puro ai fini di classificazione ospedaliera.

Test successivi svolti con un LLM dotato di un maggior numero di parametri (\texttt{nous-hermes2:34b}) hanno prodotto storie più sintetiche e rispettose delle regole sintattiche del prompt, immettendo tuttavia alterazioni concettuali sottili e altrettanto problematiche:

\begin{minted}[fontsize=\scriptsize]{json}
{
  "case_id": "5870415",
  "story_text": "The patient, an 87-year-old individual, underwent a series of 
examinations... The first event was a complete abdomen ultrasound... 
Simultaneously, a Computerized tomography scan of the skull without contrast medium was executed. 
After 728160 seconds, another complete abdomen ultrasound was conducted. 
Following this, a Computerized tomography scan of the skull with contrast medium was carried out..."
}
\end{minted}

In questo secondo test, il modello linguistico ha dedotto arbitrariamente dettagli non presenti nel database clinico, specificando la presenza o l'assenza del contrasto in una diagnostica di routine (\textit{``with contrast medium / without contrast medium''}). Benché semanticamente pertinenti in un caso cardio-vascolare, queste integrazioni scostano la traccia artificiale dall'estrazione originale, alterando il vocabolario disponibile a valle e compromettendo la confrontabilità con le tracce reali.

La convalida empirica e i successivi processi di \textit{fine-tuning} incrociato hanno sancito il fallimento dell'ipotesi predetta. L'avvalersi dello \textit{Storytelling} sintetico arricchito da tali artefatti informativi ha comportato un degrado strutturale delle metriche computazionali. La rumorosità lessicale iniettata dai Large Language Models confonde l'attenzionalità temporale del Transformer, peggiorandone la capacità selettiva e compromettendo di fatto i valori di \textit{Balanced Accuracy}. Per via di queste problematiche, l'adozione testuale definitiva in questa tesi usa esclusivamente script di parsing lineari (\texttt{story\_generator.py}), assicurando una perfetta aderenza matematica ed enumerabile al tracciato loggato originario.


% --- Appendici, Bibliografia e Ringraziamenti ---
\backmatter

% Se hai appendici, mettile qui
\appendix
% appendici.tex - Appendici della tesi

\chapter{Dettagli implementativi}

In questa appendice si forniscono dettagli aggiuntivi sull'implementazione dei modelli e degli esperimenti presentati nella tesi.

\section{Configurazione Hardware}

L'addestramento e la valutazione del modello BERT per la classificazione delle tracce cliniche, insieme alle computazionalmente costose sessioni di Explainable AI (Integrated Gradients), sono state condotte sulla seguente workstation:
\begin{itemize}
    \item CPU: AMD Ryzen 9 5900X
    \item GPU: NVIDIA GeForce RTX 4090 (24 GB VRAM)
    \item RAM: 64 GB
    \item Sistema Operativo: Linux Ubuntu 24.04.3 LTS (ambiente WSL su Windows)
\end{itemize}

L'elevata capacità di VRAM della GPU è risultata fondamentale per l'impiego del modello Transformer \texttt{bert-base-uncased} e per garantire l'efficiente esecuzione dell'algoritmo IG Adattivo fino a 5500 steps di campionamento.

\section{Configurazione Software e Linguaggi}

L'intera pipeline metodologica, dal preprocessing dei log XES fino all'analisi dei risultati XAI, è stata sviluppata in ambiente Python:
\begin{itemize}
    \item Python 3.11.13
    \item PyTorch 2.8 (come framework principale per il Deep Learning e il Fine Tuning)
    \item HuggingFace \texttt{transformers} (per caricamento e gestione del modello BERT)
    \item Captum 0.8 (libreria per il calcolo degli Integrated Gradients)
    \item Optuna 4.5 (per l'ottimizzazione bayesiana degli iperparametri)
    \item Scikit-learn 1.7.2 (per le metriche di valutazione come \textit{Balanced Accuracy} e lo Stratified K-Fold CV)
\end{itemize}

\section{Dettagli sull'Hyperparameter Tuning}

Il processo di tuning, gestito da Optuna e rifinito manualmente, aveva l'obiettivo di ottimizzare le prestazioni rispettando le architetture originali del Transformer. Di seguito sono riportati gli iperparametri finali del modello classificatore:

\begin{itemize}
    \item Base model: \texttt{bert-base-uncased} (12 layers, 768 hidden size, 12 attention heads)
    \item Batch size: Adattiva (limite VRAM)
    \item Optimizer: AdamW
    \item Funzione di costo: Focal Loss ($\gamma$ e $\alpha$ bilanciati per amplificare l'errore sulla classe $LOS \ge 20$)
    \item Validation strategy: Stratified 5-Fold Cross Validation
\end{itemize}

\chapter{Data Augmentation con LLM Generativi}

Come descritto nel Capitolo \ref{ch:metodologia}, è stato condotto un esperimento per mitigare lo sbilanciamento di classe arricchendo il dataset con "storie cliniche" sintetizzate.

\section{Pipeline dell'Esperimento}
I modelli generativi \texttt{phi-4}, \texttt{gemma-3} e \texttt{Ollama} sono stati eseguiti localmente. Il parametro di Temperatura è stato impostato prossimo allo zero ($T \approx 0.1$) per circoscrivere l'effetto stocastico ed arginare pericolose allucinazioni mediche. Alle reti veniva fornita una traccia reale estratta da un parser per il log clinico in formato XES, con l'istruzione di creare una rappresentazione in linguaggio naturale della traccia del paziente.

\section{Esito e Degradazione delle Metriche}
Nonostante la supervisione semantica, l'inserimento dei dati aumentati nel processo di training ha paradossalmente deteriorato la metrica di \textit{Balanced Accuracy}. Il calo del potere discriminatorio di BERT evidenzia come, in domini d'alta sensibilità temporale come le sequenze Process Mining ospedaliere, lievi fluttuazioni sintattiche indotte dall'LLM generativo sfalsino irreparabilmente i vettori di attenzione attesi della \textit{classification head}, dimostrando che le variazioni non rispecchiano fedelmente il pattern intrinseco (\textit{ground truth}) reale della degenza. Per questo motivo, la data augmentation sintetica è stata scartata dalla pipeline definitiva.


\bibliographystyle{unsrtnat}
\bibliography{bibliografia} 

\chapter*{Ringraziamenti}
% ringraziamenti.tex - Ringraziamenti

Desidero ringraziare il Professor Giorgio Leonardi per la guida e il supporto forniti durante tutto il percorso di questa tesi.

Un ringraziamento particolare va alla mia famiglia per il sostegno costante durante gli anni di studio.

Ringrazio inoltre i miei colleghi e compagni di corso per gli stimolanti confronti e le collaborazioni che hanno arricchito la mia esperienza universitaria.


\end{document}